\noindent\underline{Proof.}
In view of the definition (5.4.3) of the $b_j$, it is enough to prove
\begin{equation}\tag{5.4.6}
\left\{
\begin{aligned}
a_0&=1,\\
\sum_{i=0}^{a}(-1)^iS_i\,a_{j-i}&=0
&&\text{if }j>0.
\end{aligned}
\right.
\end{equation}
Clearly $a_0=1$.  For a fixed $j>0$, the second relation is obtained
by taking the coefficient of $T^{a-1}$, using (5.4.1), in the
relation
\begin{equation}\tag{5.4.7}
\sum_{i=0}^{a}(-1)^iS_iT^{a-1+j-i}=0
\qquad\text{in }B.
\end{equation}
\footnote{The print omits the upper summation bound $a$ in (5.4.7);
the relation is (5.4.2) multiplied by $T^{j-1}$.}
This relation is obtained by multiplying (5.4.2) by $T^{j-1}$,
which proves the lemma.

\subsubsection*{Corollary 5.5.}
\underline{One has}
\begin{equation}\tag{5.5.1}
\deg(\check X)
=\frac{(-1)^n}{\deg(\varphi)}
\int_X\frac{(1+H)^{r-1}}{C(N)}.
\end{equation}

\noindent\underline{Proof.}
Apply (5.3.3), with $M=N$ the normal sheaf of
$X\hookrightarrow\mathbb P^r$, to (5.2.8).

For a smooth quasi-projective variety $S$, write
\begin{equation}\tag{5.5.2}
C(S)\stackrel{\mathrm{def}}{=}C(T_S)
\end{equation}
for its total Chern class, where $T_S$ denotes the tangent bundle of
$S$.

Dualizing the exact sequence (5.1.2) gives
\begin{equation}\tag{5.5.3}
C(\mathbb P^r)=(1+H)^{r+1}
\qquad\text{in }\operatorname{CH}(\mathbb P^r).
\end{equation}
Dualizing the exact sequence (5.1.1) gives
\begin{equation}\tag{5.5.4}
C(N)\,C(X)=C(\mathbb P^r)|_X
\qquad\text{in }\operatorname{CH}(X).
\end{equation}
Combining (5.5.3) and (5.5.4), formula (5.5.1) becomes:

\subsubsection*{Corollary 5.6.}
\underline{One has}
\begin{equation}\tag{5.6.1}
\deg(\check X)
=\frac{(-1)^n}{\deg(\varphi)}
\int_X\frac{C(X)}{(1+H)^2}.
\end{equation}
\footnote{The print has $(-1)^r$.  Substitution of (5.5.3) and
(5.5.4) into (5.5.1) does not alter its sign $(-1)^n$; formula
(5.7.3) immediately below confirms the same exponent.}

\subsubsection*{5.7.}
For a smooth projective variety $S$, define its
\underline{Euler--Poincaré characteristic} by
\begin{equation}\tag{5.7.1}
\chi(S)=\int_SC(S).
\end{equation}

\subsubsection*{Proposition 5.7.2.}
\underline{Let $Y$ be a smooth hyperplane section of $X$, and let
$\Delta$ be a smooth section of $X$ by a codimension-two linear
subspace of $\mathbb P^r$.  Then}
\begin{equation}\tag{5.7.3}
\deg(\check X)
=\frac{(-1)^n}{\deg(\varphi)}
\bigl[\chi(X)+\chi(\Delta)-2\chi(Y)\bigr].
\end{equation}

